\section{Теоретическая сложность операций над матрицами и векторами}
\label{sec:TheoreticalComplexity}
\tikzsetfigurename{LinearAlgebra_Complexity_}

\mytodo{Переписать раздел. Добавить про матрицу на вектор. Отсылки к мелкозернистой сложности. Наверное, не стоит разбирать Штрассена и прочее. Просто ссылки.}

В рамках такого раздела теории сложности, как мелкозернистая сложность (fine-grained complexity) задача умножения двух матриц оказалась достаточно важной, так как через вычислительную сложность этой задачи можно оценить сложность большого класса различных задач.
С примерами таких задач можно ознакомиться в работе~\sidecite{Williams:2010:SEP:1917827.1918339}. Поэтому рассмотрим алгоритмы нахождения произведения двух матриц более подробно.
Далее для простоты мы будем предполагать, что перемножаются две квадратные матрицы одинакового размера $n \times n$.

Для начала построим наивный алгоритм, сконструированный на основе определения произведения матриц.
Такой алгоритм представлен на листинге~\ref{algo:MxM}.
\marginnote{TODO: Оформление алгоритмов точно надо обсудить, потому что я в этом мало понимаю.}
Его работу можно описать следующим образом: для каждой строки в первой матрице и для каждого столбца в второй матрице найти сумму произведений соответствующих элементов.
Данная сумма будет значением соответствующей ячейки результирующей матрицы.

%\begin{algorithm}{Наивное умножение матриц}{MxM}
%    \begin{pseudo}[]
%        \kw{function} \pr{MatrixMult}(M_1, M_2, G = (S, \oplus, \otimes)) \\+
%        $M_3 = $ пустая матрица размера $n \times n$ \\
%        \kw{for} $i \in [0 \rng n - 1]$ \\+
%        \kw{for} $j \in [0 \rng n - 1]$ \\+
%        \kw{for} $k \in [0 \rng n - 1]$ \\+
%        $M_3[i, j] = M_3[i, j] \oplus (M_1[i, k] \otimes M_2[k, j])$ \\---
%        \kw{return} $M_3$
%    \end{pseudo}
%\end{algorithm}

Сложность наивного произведения двух матриц составляет $O(n^3)$ из-за тройного вложенного цикла, где каждый уровень вложенности привносит $n$ итераций.
Но можно ли улучшить этот алгоритм?
Первый положительный ответ был опубликован Ф. Штрассен в 1969 году~\sidecite{Strassen1969}.
Сложность предложенного им алгоритма~--- $O(n^{\log_2 7}) \approx O(n^{2.81})$.
Основная идея~--- рекурсивное разбиение исходных матриц на блоки и вычисление их произведения с помощью только 7 умножений, а не 8.

Рассмотрим алгоритм Штрассена более подробно.
Пусть $A$ и $B$~--- две квадратные матрицы размера $2^n \times 2^n$ над кольцом $R=(S, \oplus, \otimes)$.
Если размер умножаемых матриц не является натуральной степенью двойки, то дополняем исходные матрицы дополнительными нулевыми строками и столбцами.
Наша задача найти матрицу $C = A \cdot B$.

Разделим матрицы $A, B$ и $C$ на четыре равные по размеру блока.
\[
    A =
    \begin{pmatrix}
        A_{1,1} & A_{1,2} \\
        A_{2,1} & A_{2,2}
    \end{pmatrix},
    \quad
    B =
    \begin{pmatrix}
        B_{1,1} & B_{1,2} \\
        B_{2,1} & B_{2,2}
    \end{pmatrix},
    \quad
    C =
    \begin{pmatrix}
        C_{1,1} & C_{1,2} \\
        C_{2,1} & C_{2,2}
    \end{pmatrix}
\]

По определению произведения матриц выполняются следующие равенства.
\marginnote{TODO: Вообще можно попробовать раскидать на 2 столбца}
\begin{align*}
    C_{1, 1} & = A_{1, 1} \cdot B_{1, 1} + A_{1, 2} \cdot B_{2, 1} \\
    C_{1, 2} & = A_{1, 1} \cdot B_{1, 2} + A_{1, 2} \cdot B_{2, 2} \\
    C_{2, 1} & = A_{2, 1} \cdot B_{1, 1} + A_{2, 2} \cdot B_{2, 1} \\
    C_{2, 2} & = A_{2, 1} \cdot B_{1, 2} + A_{2, 2} \cdot B_{2, 2}
\end{align*}

Данная процедура не даёт нам ничего нового с точки зрения вычислительной сложности.
Но мы можем двинуться дальше и определить следующие элементы.
\begin{align*}
    P_1 & \equiv (A_{1, 1} + A_{2, 2}) \cdot (B_{1, 1} + B_{2, 2}) \\
    P_2 & \equiv (A_{2, 1} + A_{2, 2}) \cdot B_{1, 1}              \\
    P_3 & \equiv A_{1, 1} \cdot (B_{1, 2} - B_{2, 2})              \\
    P_4 & \equiv A_{2, 2} \cdot (B_{2, 1} - B_{1, 1})              \\
    P_5 & \equiv (A_{1, 1} + A_{1, 2}) \cdot B_{2, 2}              \\
    P_6 & \equiv (A_{2, 1} - A_{1, 1}) \cdot (B_{1, 1} + B_{1, 2}) \\
    P_7 & \equiv (A_{1, 2} - A_{2, 2}) \cdot (B_{2, 1} + B_{2, 2})
\end{align*}

Используя эти элементы мы можем выразить блоки результирующей матрицы следующим образом.
\begin{align*}
    C_{1, 1} & = P_1 + P_4 - P_5 + P_7 \\
    C_{1, 2} & = P_3 + P_5             \\
    C_{2, 1} & = P_2 + P_4             \\
    C_{2, 2} & = P_1 - P_2 + P_3 + P_6
\end{align*}

При таком способе вычисления мы получаем на одно умножение подматриц меньше, чем при наивном подходе.
Это и приводит, в конечном итоге, к улучшению сложности всего алгоритма, который основывается на рекурсивном повторении проделанной выше процедуры.

\marginnote{TODO: здесь \textbackslash{}sidecite не влезает}
Впоследствии сложность постепенно понижалась в ряде работ, таких как~\cite{Pan1978,BiniCapoRoma1979,Schonhage1981,CoppWino1982,CoppWino1990}.
Было введено специальное обозначение для показателя степени в данной оценке: $\omega$.
То есть сложность умножения матриц~--- это $O(n^\omega)$, и задача сводится к уменьшению значения $\omega$.
В настоящее время работа над уменьшением показателя степени продолжается и сейчас уже предложены решения с $\omega < 2.373$%
\sidenote{
    В данной области достаточно регулярно появляются новые результаты, дающие сравнительно небольшие, в терминах абсолютных величин, изменения.
    Так, в 2021 была представлена работа, улучшающая значение $\omega$ в пятом знаке после запятой~\cite{alman2020refined}.
    Несмотря на кажущуюся несерьёзность результата, подобные работы имеют большое теоретическое значение, так как улучшают наше понимание исходной задачи и её свойств.}%
.

Всё тем же Ф. Штрассеном ещё в 1969 году была выдвинута гипотеза о том, что для достаточно больших $n$ существует алгоритм, который для любого сколь угодно маленького наперёд заданного $\varepsilon$ перемножает матрицы за $O(n^{2+\varepsilon})$.
На текущий момент ни доказательства, ни опровержения этой гипотезы не предъявлено.

Важной особенностью указанного выше направления улучшения алгоритмов является то, что оно допускает использование (и даже основывается на использовании) более богатых алгебраических структур, чем требуется для определения умножения двух матриц.
Так, уже алгоритм Штрассена использует операцию вычитания, что приводит к необходимости иметь обратные элементы по сложению, а значит определять матрицы над кольцом.
Хотя для исходного определения (\ref{def:MxM}) достаточно более бедной структуры.
При этом, часто, структуры, возникающие в прикладных задачах кольцами не являются.
\marginnote{TODO: Не кажется ли что текста на полях слишком много и его можно прямо в главу вписать?}
Примерами могут служить тропическое (или $\{min, +\}$) полукольцо, играющее ключевую роль в тропической математике, или булево ($\{\lor, \land\}$) полукольцо, возникающее, например, при работе с отношениями%
\sidenote{
    Вообще говоря, в некоторых прикладных задачах возникают структуры, не являющиеся даже полукольцом.
    Предположим, что есть три различных множества $S_1$, $S_2$ и $S_3$ и две двухместные функции $f: S_1 \times S_2 \to S_3$ и $g: S_3 \times S_3 \to S_3$.
    Этого достаточно, чтобы определить произведение двух матриц $M_1$ и $M_2$, построенных из элементов множеств $S_1$ и $S_2$ соответственно.
    Результирующая матрица будет состоять из элементов $S_3$.
    Как видно, функции не являются бинарными операциями в смысле нашего определения.
    Несмотря на кажущуюся экзотичность, подобные структуры возникают на практике при работе с графами и учитываются, например, в стандарте GraphBLAS (\url{https://graphblas.github.io/}), где, кстати, называются полукольцами, что выглядит не вполне корректно.}%
.
Значит, описанные выше решения не применимы и вопрос о существовании алгоритма с менее чем кубической сложностью снова актуален.

В попытках ответить на этот вопрос появились так называемые комбинаторные алгоритмы умножения матриц%
\sidenote{
    В противовес описанным выше, не являющимся комбинаторными.
    Стоит отметить, что строгое определение комбинаторных алгоритмов отсутствует, хотя этот термин и получил широкое употребление.
    В частности, Н.~Бансал (Nikhil Bansal) и Р.~Уильямс (Ryan Williams) в работе~\cite{5438580} дают определение комбинаторного алгоритма, но тут же замечают следующее: \enquote{We would like to give a definition of \enquote{combinatorial algorithm}, but this appears elusive. Although the term has been used in many of the cited references, nothing in the literature resembles a definition. For the purposes of this paper, let us think of a \enquote{combinatorial algorithm} simply as one that does not call an oracle for ring matrix multiplication.}.
    Ещё один вариант определения и его обсуждение можно найти в~\cite{das2018lower}.}%
.
Классический результат в данной области~--- это алгоритм четырёх русских, предложенный В. Л. Арлазаровым, Е. А. Диницем, М. А. Кронродом и И. А. Фараджевым в 1970 году~\cite{ArlDinKro70}, позволяющий перемножить матрицы над конечным полукольцом за $O(n^3/\log n)$.
Лучшим результатом%
\sidenote{
В работе~\cite{das2018lower} предложен алгоритм со сложностью $\Omega(n^{7/3}/2^{O(\sqrt{\log n})})$, однако авторы утверждают, что сами не уверены в комбинаторности предложенного решения.
По-видимому, полученные результаты ещё должны быть проверены сообществом.}
в настоящее время является алгоритм со сложностью%
\sidenote{Нотация $\hat{O}$ скрывает $poly(\log\log)$ коэффициенты.} $\hat{O}(n^3/\log^4 n)$~\cite{10.1007/978-3-662-47672-7_89}.

Как видим, особенности алгебраических структур накладывают серьёзные ограничения на возможности конструирования алгоритмов.
Отметим, что, хотя, в указанных случаях и предлагаются решения лучшие, чем наивное кубическое, они обладают принципиально разной асимптотической сложностью.
В первом случае сложность оценивается полиномом, степень которого меньше третьей.
Такие решения принято называть \emph{истинно субкубическими} (truly subcubic).
В то время как в случае комбинаторных алгоритмов степень полинома остается прежней, третьей, хотя сложность и уменьшается на логарифмический фактор.
Такие решения принято называть \emph{слегка субкубическими} (mildly subcubic).
Естественный вопрос о существовании истинно субкубического алгоритма перемножения матриц над полукольцами (или же комбинаторного перемножения матриц) всё ещё не решён%
\sidenote{Один из кандидатов~--- работа~\cite{das2018lower}, однако на текущий момент предложенное в ней решение  требует проверки.}.
